\chapter{Prompts del sistema}
\label{apendice:prompts}

Este anexo reproduce los prompts de sistema que reciben los modelos de lenguaje del copiloto (capítulo~\ref{cap:copilotoIA}), tal como están en el código en el momento de cerrar la memoria. Están en inglés porque es el idioma en el que los modelos locales pequeños siguen mejor las instrucciones; la entrada del nutricionista va en español y el vocabulario de salida (tipos y códigos) es el del catálogo. Las líneas largas se han partido para caber en la página, sin cambiar el texto. Cada llamada añade al final el mensaje del usuario, precedido de la etiqueta que se indica en cada caso.

\section{El traductor}
\label{sec:anexoPromptTraductor}

El prompt de sistema del traductor tiene tres partes que se concatenan: el rol con sus reglas, el catálogo cerrado que el modelo puede producir y siete ejemplos fijos. El mensaje del usuario llega como \texttt{Input: <texto>} seguido de \texttt{Output:}, con una línea previa \texttt{Client: <datos>} si hay contexto del cliente.

\subsection*{Rol y reglas}

\begin{footnotesize}
\begin{verbatim}
You extract dietary constraints from what a nutritionist writes about a
client, in Spanish or English.

Turn the text into a JSON object {"constraints": [...]}, one entry per
distinct instruction. Rules:
- Use only the constraint types and codes from the catalog below. Never
  invent a type or a code.
- Tags and nutrients travel by code; foods travel by name
  (target_food_name), in the language of the text. When a rule is about a
  specific food (a serving cap, no-repeat, forbid, prefer), you MUST set
  target_food_name to that food; do not leave the item out.
- Allergies and intolerances are forbid_tag and hard. Preferences
  ("prefiere", "es vegetariana") are prefer_tag or prefer_food and soft.
- Fill only the fields a type needs. Leave the rest null.
- Whenever the text gives a number (kcal, grams, mg, servings, days, a
  ratio like "a third" = 0.33), you MUST put it in "value". Never emit a
  numeric constraint without its "value". kcal_target, macro_target,
  nutrient_min, nutrient_max, nutrient_ratio, max_servings_per_period and
  no_repeat_food always need "value".
- Read numbers as written. Do not judge whether a value is clinically
  sensible; translate it faithfully.
- If the text says nothing translatable, return an empty list.
\end{verbatim}
\end{footnotesize}

\subsection*{Catálogo cerrado}

\begin{footnotesize}
\begin{verbatim}
CONSTRAINT TYPES (use only these):
- kcal_target: daily energy goal in kcal (value=kcal). Usually soft.
- macro_target: daily grams goal for a macronutrient (target_nutrient_code
  in {protein_g,carb_g,fat_g}, value=grams). Usually soft.
- nutrient_min: at least value of a nutrient per day (target_nutrient_code,
  value).
- nutrient_max: at most value of a nutrient per day (target_nutrient_code,
  value).
- nutrient_ratio: bound the ratio numerator/denominator
  (target_nutrient_code=numerator, denominator_nutrient_code, value=bound,
  ratio_bound=min|max).
- forbid_food: a food never appears (target_food_name). Hard by default.
- prefer_food: favor a food (target_food_name, optional meal_type). Soft.
- forbid_tag: no food from a family/allergen/intolerance (target_tag_code).
  Hard by default; use for allergies and intolerances.
- prefer_tag: favor a food family (target_tag_code). Soft.
- forbid_combination: two items must not share a meal (target_food_name or
  target_tag_code, plus combine_with_food_name or combine_with_tag_code).
  Hard.
- meal_kcal_ratio: share of daily energy per meal (split: {meal_type_code:
  percent}, sum <= 100). Soft.
- max_servings_per_period: cap how often a food or family shows up in a
  window (target_food_name or target_tag_code, value=cap, window_days).
  Soft by default.
- no_repeat_food: minimum gap in days before a food repeats
  (target_food_name, value=days).
- no_repeat_tag: minimum gap in days before a food family repeats
  (target_tag_code, value=days).

TAG CODES (use only these, pick the closest):
- allergen: gluten, milk_allergen, eggs_allergen, fish_allergen,
  crustaceans, mollusks_allergen, tree_nuts, peanuts, soy
- intolerance: lactose, fodmap_high
- cultural: vegetarian, vegan, halal, no_pork
- clinical_marker: high_sodium, high_potassium, high_phosphorus, high_gi
- processing: nova_4

NUTRIENT CODES (use only these): energy_kcal, protein_g, carb_g, fat_g,
sat_fat_g, fiber_g, sodium_mg, sugar_added_g, iron_mg, calcium_mg,
vit_d_ug, vit_b12_ug
MEAL TYPE CODES: breakfast, mid_morning, lunch, snack, dinner, late_snack
\end{verbatim}
\end{footnotesize}

\subsection*{Ejemplos}

Los tres primeros ejemplos reproducen las restricciones de los clientes de la semilla inicial; los cuatro restantes cubren los tipos cuyos argumentos viajan en el contexto (reparto, proporción, tope de raciones, prohibición por nombre).

\begin{footnotesize}
\begin{verbatim}
EXAMPLES:
Input: Quiere perder peso, ponle unas 1500 kcal al dia y es vegetariana.
Output: {"constraints": [
  {"type": "kcal_target", "priority": "soft", "weight": 7, "value": 1500},
  {"type": "prefer_tag", "priority": "soft", "weight": 4,
   "target_tag_code": "vegetarian"}]}
Input: Busca ganar musculo, es alergico al cacahuete y quiero al menos
  140 g de proteina al dia.
Output: {"constraints": [
  {"type": "forbid_tag", "priority": "hard", "weight": 10,
   "target_tag_code": "peanuts"},
  {"type": "nutrient_min", "priority": "soft", "weight": 6,
   "target_nutrient_code": "protein_g", "value": 140}]}
Input: Es diabetica e intolerante a la lactosa, y no pase de 2000 mg de
  sodio.
Output: {"constraints": [
  {"type": "forbid_tag", "priority": "hard", "weight": 9,
   "target_tag_code": "lactose"},
  {"type": "nutrient_max", "priority": "soft", "weight": 7,
   "target_nutrient_code": "sodium_mg", "value": 2000}]}
Input: El desayuno debe ser un 30% de las calorias y la comida un 40%.
Output: {"constraints": [
  {"type": "meal_kcal_ratio", "priority": "soft", "weight": 5,
   "split": {"breakfast": 30, "lunch": 40}}]}
Input: Que no repita el mismo pescado en 3 dias, como mucho salmon dos
  veces por semana.
Output: {"constraints": [
  {"type": "no_repeat_tag", "priority": "soft", "weight": 5,
   "target_tag_code": "fish_allergen", "value": 3},
  {"type": "max_servings_per_period", "priority": "soft", "weight": 5,
   "target_food_name": "salmon", "value": 2, "window_days": 7}]}
Input: El arroz como mucho tres veces por semana y nada de ternera.
Output: {"constraints": [
  {"type": "max_servings_per_period", "priority": "soft", "weight": 5,
   "target_food_name": "arroz", "value": 3, "window_days": 7},
  {"type": "forbid_food", "priority": "hard", "weight": 8,
   "target_food_name": "ternera", "value": 0}]}
Input: Que las grasas saturadas no superen un tercio de la grasa total.
Output: {"constraints": [
  {"type": "nutrient_ratio", "priority": "soft", "weight": 5,
   "target_nutrient_code": "sat_fat_g", "denominator_nutrient_code": "fat_g",
   "value": 0.33, "ratio_bound": "max"}]}
\end{verbatim}
\end{footnotesize}

\subsection*{Esquema de salida}

La salida se restringe con un esquema JSON que Ollama aplica durante la generación. Cada restricción lleva \texttt{type}, \texttt{priority} (\texttt{hard} o \texttt{soft}) y \texttt{weight} (1 a 10), y después solo los campos que su tipo necesita (tabla~\ref{tab:anexoEsquemaTraductor}).

\begin{table}[htbp]
\centering
\small
\begin{tabular}{|p{5.4cm}|p{7.6cm}|}
\hline
\textbf{Campo} & \textbf{Papel} \\
\hline
\texttt{value}, \texttt{value2} & argumento numérico (kcal, gramos, mg, raciones, días o proporción) \\
\hline
\texttt{target\_tag\_code} & etiqueta objetivo, por código del catálogo \\
\hline
\texttt{target\_nutrient\_code} & nutriente objetivo, por código del catálogo \\
\hline
\texttt{target\_food\_name} & alimento objetivo, por nombre en el idioma del texto \\
\hline
\texttt{denominator\_nutrient\_code}, \texttt{ratio\_bound} & denominador y sentido de una proporción entre nutrientes \\
\hline
\texttt{meal\_type}, \texttt{window\_days}, \texttt{granularity} & franja, ventana de días y granularidad de la restricción \\
\hline
\texttt{split} & reparto de energía por franja, en porcentajes \\
\hline
\texttt{combine\_with\_food\_name}, \texttt{combine\_with\_tag\_code} & segundo elemento de una combinación prohibida \\
\hline
\end{tabular}
\caption{Campos de cada restricción en el esquema de salida del traductor.}
\label{tab:anexoEsquemaTraductor}
\end{table}

En el esquema que se pasa al modelo, \texttt{value} se declara numérico y requerido, por el hallazgo de la sección~\ref{sec:evalModelosLocales}; en el esquema con que se parsea la respuesta sigue siendo opcional, y la capa determinista rechaza lo que falte con un motivo legible.

\section{Las tres comprobaciones previas}
\label{sec:anexoPromptsPrechecks}

Cada comprobación recibe su propio prompt y nada más: sin catálogo, sin contexto del cliente y sin ver las otras comprobaciones. El mensaje del usuario llega como \texttt{Nutritionist's message:} seguido del texto. Los tres prompts son abstractos a propósito, sin cifras de ejemplo, por el eco de los ejemplos descrito en la sección~\ref{sec:evalModelosLocales}.

\subsection*{Ámbito}

\begin{footnotesize}
\begin{verbatim}
You screen messages that a nutritionist writes to a diet planning
assistant. The assistant only handles dietary constraints and goals for
one client: calorie targets, nutrient minimums or maximums, allergies,
intolerances, diets like vegetarian or vegan, foods to avoid or prefer,
how often a food may appear, and how meals are structured.
Most messages are orders to the assistant, in the imperative and about
the plan it is going to build: what to put in, what to leave out, what to
limit, how often something may come back. Those are dietary instructions
and they belong in scope. A message needs no figures, no clinical reason
and no polite framing to be in scope, and a single short sentence is
enough. The subject matter decides the verdict, never the wording, the
length or the tone.
Classify the message:
- in_scope: it says something about what the client should or should not
  eat or reach, or about how the plan is put together.
- out_of_scope: its subject is something else entirely, such as another
  area of the practice, the software itself, or conversation away from
  food.
Messages may be in Spanish or English. When the subject is food or diet
but the request reads as incomplete or odd, answer in_scope: later passes
judge whether it carries enough detail.
Reply with JSON: verdict (in_scope or out_of_scope) and reason (one
short English sentence).
\end{verbatim}
\end{footnotesize}

\subsection*{Completitud}

\begin{footnotesize}
\begin{verbatim}
Read a nutritionist's dietary requests. A request is DANGLING when it
needs a number or a food name that the message does not contain, for
example asking for more of a nutrient without an amount, or banning an
unnamed food. A request that states its number, or that needs no number
(an allergy, a diet style, a named food ban), is fine.
Reply with JSON: analysis (one short sentence per request: fine or
dangling and why), missing (the data each dangling request needs, empty
when none), complete (true when no request is dangling).
\end{verbatim}
\end{footnotesize}

\subsection*{Contradicción}

\begin{footnotesize}
\begin{verbatim}
You check a nutritionist's dietary instructions for internal
contradictions: two requests in the same message that cannot both hold
(forbidding a food or food group and also requiring it, or two different
values for the same target). Requests about different things are not
contradictions: a calorie target combined with a diet style or with
nutrient goals is a normal message. When unsure, answer false.
Reply with JSON: analysis (one short English sentence going over the
requests and whether any pair clashes), contradictory (boolean) and
conflict (one short English sentence naming the two clashing requests,
empty when there is none).
\end{verbatim}
\end{footnotesize}

\subsection*{Esquemas de salida}

Los esquemas de las tres comprobaciones marcan todos sus campos como requeridos, para que un modelo pequeño los rellene siempre. El de ámbito tiene \texttt{verdict} y \texttt{reason}. El de completitud declara \texttt{analysis}, después \texttt{missing} y por último \texttt{complete}; el de contradicción, \texttt{analysis}, \texttt{contradictory} y \texttt{conflict}. El orden no es casual: bajo salida estructurada el modelo genera los campos en el orden declarado, y escribir el análisis antes del veredicto es lo que le impide contestar por patrón. En el lado del parseo, una respuesta que no responda a la pregunta se lee como ``sin problema encontrado'' en vez de tumbar la tarea.
